\section{Outcome Verification: partiendo de GSM8K}

\subsection{Training Verifiers to Solve Math Word Problems}

El trabajo de Cobbe et al. aporta a la vez el conjunto de datos GSM8K y un método sistemático para entrenar un verifier. GSM8K contiene alrededor de 8.5K problemas de matemáticas de primaria en forma de texto, cada uno acompañado de un proceso de razonamiento en lenguaje natural y una respuesta numérica final.

\videofigure{figures/gsm8k-contributions.jpg}{GSM8K y el outcome verifier constituyen las dos contribuciones centrales de este trabajo.}{00:02:10--00:04:26}

El flujo central es: primero se entrena el generator; para cada problema de entrenamiento se muestrean numerosas completions; se asigna una etiqueta binaria según si la respuesta final coincide con el ground truth; después se entrena el verifier para predecir si toda la solución es correcta.

\videofigure{figures/generator-verifier-pipeline.jpg}{El generator produce candidatos etiquetados, y el verifier aprende a puntuar la solución completa.}{00:04:26--00:06:20}

\subsection{El objetivo conjunto del Verifier}

El verifier puede agregar una cabeza de clasificación sobre el backbone del modelo de lenguaje, conservando al mismo tiempo el objetivo de next-token:

$$
\mathcal L(\phi)
=\mathcal L_{\mathrm{LM}}(\phi)
+\lambda\,\mathcal L_{\mathrm{cls}}(\phi),
$$

$$
\mathcal L_{\mathrm{cls}}
=-z\log v_{\phi}(x,y)-(1-z)\log(1-v_{\phi}(x,y)).
$$

\begin{itemize}
    \item $\phi$: parámetros del verifier;
    \item $z\in\{0,1\}$: si la respuesta final es correcta;
    \item $v_{\phi}(x,y)$: probabilidad de que el candidato sea correcto;
    \item $\lambda$: peso entre el modelado de lenguaje y la clasificación de corrección.
\end{itemize}

\videofigure{figures/verifier-objective.jpg}{El verifier modela a la vez el texto y la corrección de la solución completa.}{00:06:20--00:08:16}

\begin{knowledgebox}{Por qué conservar el objetivo de modelado de lenguaje}
Una clasificación pura podría aprender únicamente atajos como el formato, la longitud o palabras clave de la respuesta final. El objetivo de modelado de lenguaje obliga a la representación a seguir entendiendo el texto completo del razonamiento, y además funciona como regularización cuando los datos son limitados. Sin embargo, no elimina automáticamente los shortcuts; sigue siendo necesario revisarlos mediante ejemplos adversarios y evaluación fuera de distribución.
\end{knowledgebox}

\subsection{Los datos de entrenamiento deben provenir de la propia distribución de errores del modelo}

Si el verifier solo observa errores obvios construidos manualmente, no necesariamente podrá reconocer los patrones de fallo reales del generator. Este método muestrea numerosas completions del modelo para cada problema, de modo que el conjunto de entrenamiento incluya negativos difíciles que "se parecen mucho a una solución correcta".

\begin{warningbox}{El Verifier no puede limitarse a obtener buen desempeño en datos estáticos}
Actualizar el generator, cambiar el prompt o aumentar la longitud del razonamiento modifican la distribución de errores. Un verifier preciso sobre candidatos antiguos puede fallar con un modelo nuevo. Por ello, los datos del verifier deben muestrearse continuamente del generator actual, y su evaluación debe llevar control de versiones.
\end{warningbox}

\subsection{Best-of-N y cómputo en tiempo de inferencia}

En el momento de la inferencia se generan $N$ candidatos y se elige aquel con la puntuación más alta del verifier:

$$
y^*=\arg\max_{y_j\sim\pi_\theta(\cdot\mid x)} v_\phi(x,y_j).
$$

A medida que $N$ aumenta, mejora la cobertura de candidatos; mientras el verifier conserve capacidad de discriminación en su ordenamiento, la precisión final también mejora. Se ha observado experimentalmente que el desempeño sigue mejorando hasta con cientos de completions.

\videofigure{figures/testtime-compute.jpg}{Aumentar el número de candidatos sigue mejorando la precisión tras la selección del verifier.}{00:16:50--00:19:04}

\subsection{La limitación fundamental de outcome supervision}

La etiqueta de la respuesta final no puede indicar en qué paso se produjo el error del razonamiento. Una trayectoria puede estar completamente equivocada al inicio y aun así llegar por coincidencia al valor numérico correcto; o bien el proceso puede ser esencialmente correcto y fallar solo en un error aritmético al final. El ORM etiqueta ambos casos como positivo y negativo respectivamente, lo que produce una asignación de crédito muy tosca.

\subsection{Resumen de la sección}

El outcome verifier demostró por primera vez que learned verification puede convertir el inference sampling en una mejora real de desempeño, y constituye la base de la recompensa de RL en trabajos posteriores. Pero mirar solo la respuesta final confunde el proceso correcto, la corrección accidental y el error local, lo cual impulsó a la investigación a orientarse hacia process supervision.
